\section{Model and the location degeneracy}\label{sec:model}

\subsection{Mixed models and joint models}\label{sec:setup}

Index subjects by $i = 1, \dots, N$ and subject $i$'s longitudinal
observations by $j = 1, \dots, n_i$, taken at times $t_{ij}$. Subject $i$
has a fixed-effects design $X_i \in \R^{n_i \times p}$ and a random-effects
design $Z_i \in \R^{n_i \times q}$, and the linear predictor is
\begin{equation}\label{eq:lp}
\eta_{ij} = x_{ij}^{\top}\beta + z_{ij}^{\top} b_i ,
\qquad b_i \sim \N_q(0, \Sigma), \quad \Sigma = L L^{\top},
\end{equation}
with $L = \operatorname{diag}(\sigma_b)\,L_{\mathrm{corr}}$ lower
triangular; for a random intercept and slope, $\sigma_b = (\sigma_0,
\sigma_1)$ and $\rho$ is their correlation. The response is Gaussian with mean $\eta_{ij}$ and variance
$\sigma_e^2$, or follows a generalized linear model with linear predictor
$\eta_{ij}$. Collect the random effects as the rows of $B \in \R^{N \times q}$.
The non-centred parameterization samples
\begin{equation}\label{eq:ncp}
\Bstd \in \R^{N \times q}, \quad \text{entries iid } \N(0,1),
\qquad B = \Bstd L^{\top},
\end{equation}
and $\theta$ collects every other parameter, including $\beta$,
$\sigma_b$ and $L_{\mathrm{corr}}$.

In a joint longitudinal-survival model the longitudinal submodel is
Gaussian, and subject $i$'s current underlying marker value
\begin{equation}\label{eq:traj}
m_i(t) = x_i(t)^{\top}\beta + z_i(t)^{\top} b_i
\end{equation}
enters a proportional hazards model,
\begin{equation}\label{eq:hazard}
h_i(t) = h_0(t)\exp\{\gamma^{\top} w_i + \alpha\, m_i(t)\},
\end{equation}
with baseline covariates $w_i$ and association parameter $\alpha$. Here
$x_i(t)$ and $z_i(t)$ are the longitudinal design rows evaluated at time
$t$ with subject $i$'s baseline covariates. In \texttt{jmjax} the log baseline
hazard is a cubic B-spline with a second-order random-walk smoothing prior,
and the cumulative hazard is computed by Gauss--Legendre quadrature, so the
random effects reach the likelihood through $\eta_{ij}$ and through
$m_i(t)$ at each event time and each quadrature node. The priors on
$\beta$, $\gamma$ and $\alpha$ are normal, the random-effect standard
deviations have gamma or half-normal priors, and $L_{\mathrm{corr}}$ an LKJ
prior; none of what follows depends on these choices beyond the prior on
$\Bstd$ being iid standard normal.

\subsection{The degeneracy}\label{sec:degeneracy}

Take a random intercept and slope, $z_{ij} = (1, t_{ij})^{\top}$, and a
fixed-effects design with an intercept, a time column and subject-constant
covariates such as age and sex. Three kinds of shift leave every
$\eta_{ij}$ and every $m_i(t)$ unchanged:
\begin{itemize}
\item $b_{i0} \mapsto b_{i0} + c$ for all $i$, with $\beta_0 \mapsto
\beta_0 - c$;
\item $b_{i0} \mapsto b_{i0} + c\,\mathrm{age}_i$ for all $i$, with
$\beta_{\mathrm{age}} \mapsto \beta_{\mathrm{age}} - c$, and likewise for
every subject-constant covariate;
\item $b_{i1} \mapsto b_{i1} + c$ for all $i$, with $\beta_{\mathrm{time}}
\mapsto \beta_{\mathrm{time}} - c$.
\end{itemize}
The second kind adds one direction for each subject-constant covariate.
The third is present whenever there is a random slope, whatever the
covariates.
Along each shift the likelihood is exactly flat, and only the prior on the
random effects locates the split, at scale $\sigma_0/\sqrt{N}$ for the
mean intercept direction.

In the coordinates HMC samples, a shift of this kind moves all $N$
entries of one or both columns of $\Bstd$ together (for the mean intercept
direction with uncorrelated random effects, each by $c/\sigma_0$), against
one or a few coordinates of $\beta$. A diagonal mass matrix, which
is what NUTS adapts by default, can only rescale coordinates one at a time
and cannot align with such a direction. In a fit of the \texttt{pbc2} data
(Section~\ref{sec:realdata}) with an intercept and one subject-constant
covariate, the posterior correlation between $\beta_0$ and the mean random
intercept was $-0.971$, and between the covariate's coefficient and the
matching covariate-weighted combination $-0.972$. The two affected
coefficients had minimum
effective sample sizes of 441 and 526 from 4000 draws, against 2893 and
more than 4000 for the identified sums. Stan's user's guide documents the
same mechanism for a model with two intercepts under weak priors
\citep{stanug}.

The degeneracy is not a funnel. The funnel is a problem of curvature that
changes with a scale parameter; here the curvature along the ridge is
nearly constant and the problem is that the ridge is narrow in a direction
the metric cannot see. It is also not the non-identifiability of an
improper model: the prior on the random effects makes the posterior proper
and the split identified, weakly.
